% =============================================================================
% Appendix
% =============================================================================

\appendix

\section{Judge Prompt Template}
\label{app:judge-prompt}

The Sonnet-4.6 judge in \S\ref{sec:eval-judge} is invoked with a
fixed system prompt and a per-pair user message whose A/B ordering is
randomized per call. The full templates appear below verbatim;
the implementation lives in
\texttt{src/skillcacher/judge/prompt.py}.

\paragraph{System prompt.}
\begin{quote}\footnotesize\itshape
You are evaluating two outputs from an AI coding agent. The agent is
operating in an interactive multi-turn task where it uses tool calls
(JSON function-call blocks) to read files, edit code, and make
progress on software-engineering tasks.

Both outputs are responses from the SAME agent context to the SAME
prompt --- the only difference is the engine's KV-cache strategy.
Your job is to identify which output is a better next step for the
agent's task.

When comparing:
\begin{itemize}
\item A well-formed tool call (JSON specifying a function name +
  arguments) is generally better than natural-language prose,
  because the agent's task requires tool calls to make progress.
\item When both outputs are tool calls, prefer the one with more
  appropriate arguments --- sensible file paths, reasonable command
  flags, useful descriptions.
\item When both outputs are prose, prefer the one that more directly
  addresses the task.
\item If the outputs are functionally equivalent (e.g., same tool
  with cosmetically different argument wording), respond
  \texttt{EQUIVALENT}.
\end{itemize}

Respond with EXACTLY one of these tokens on the first line:
\texttt{PREFER\_A} \textbar{} \texttt{PREFER\_B} \textbar{} \texttt{EQUIVALENT}.
Then on a new line, give a 1-sentence rationale ($\leq$ 30 words).
\end{quote}

\paragraph{User template.} Rendered per pair with the task prompt
and both outputs substituted:
\begin{lstlisting}[basicstyle=\ttfamily\footnotesize, frame=single,
                   framesep=4pt, breaklines=true]
TASK PROMPT (what the agent was asked to do):
{task_prompt}

OUTPUT A:
{output_a}

OUTPUT B:
{output_b}
\end{lstlisting}

\paragraph{Parsing.} Responses are parsed with the line-anchored
regex \texttt{\^{}\textbackslash{}s*(PREFER\_A|PREFER\_B|EQUIVALENT)\textbackslash{}b}
applied in multiline mode. The label is mapped back to
\{cacheblend, no\_cache, equivalent\} via the recorded
\texttt{position\_a} for that call; anything that does not match the
regex on either of the first two lines is recorded as
\texttt{UNPARSEABLE} (zero occurrences across the 19-pair run).
Rationales are extracted as the remainder of the response after the
label token, trimmed and truncated at the first newline.

\paragraph{Calling conventions.} The judge is invoked via the
Anthropic Messages API with top-level \texttt{cache\_control:
\{type: "ephemeral"\}} on the static system prompt; per-pair user
messages bypass the cache. Position assignment is seeded
per pair (\texttt{random.Random(42 \^{} hash((capture\_id,
turn\_index)))}) so that re-running the judge call deterministically
reproduces both the A/B ordering and any cache hits on the system
prompt.

\section{Statistical Span Miner}
\label{app:miner}

The miner extracts repeated token sequences from the
\dataset{} corpus for the pre-seed dispatcher
(\S\ref{sec:miner}). Inputs: a list of token streams (one per
captured request). Parameters: a length floor
$\ell_{\min}$=256~tokens (matching
\texttt{lmcache.blend\_min\_tokens}), a frequency floor
$f_{\min}$=3~distinct streams, a MinHash Jaccard cutoff $J$=0.8,
and a shingle width $k$=5. The output is a list of
\texttt{MinedSpan} records, sorted by
$(\text{frequency descending}, \text{length descending})$, each
carrying its token ids, originating stream ids, and frequency
count. The reference implementation
(\texttt{src/\allowbreak{}skillcacher/\allowbreak{}tagging/\allowbreak{}statistical\_\allowbreak{}miner.py},
$\sim$340 LoC) follows the procedure below.

\begin{algorithm}[h]
\small
\textbf{Inputs:} streams $S_1, \dots, S_m$;
                 $\ell_{\min}$, $f_{\min}$, $J$, perms $P$, shingle width $k$.\\
\textbf{Output:} sorted list of mined spans.

\begin{enumerate}
  \item[1.] \emph{Concatenate.} Build $T = S_1 \cdot \sigma_1 \cdot S_2 \cdot \sigma_2 \cdots S_m$
    where each $\sigma_i$ is a unique negative sentinel that prevents
    suffixes from crossing stream boundaries. Maintain a parallel
    array \texttt{stream\_of\_pos} mapping each position to its
    originating stream id.

  \item[2.] \emph{Suffix array via prefix doubling.} Compute
    $\mathrm{SA}[0\,..\,n{-}1]$ such that $\mathrm{SA}[i]$ is the
    starting index of the $i$-th lex-smallest suffix of $T$. Each
    round sorts on the composite key
    $(\mathrm{rank}[i], \mathrm{rank}[i{+}k])$, doubling $k$ each
    iteration; terminates when all ranks are unique. Total work
    $O(n \log^2 n)$, memory $O(n)$. (The naive sort by full suffix
    string is $O(n^2)$ in expectation on this corpus and was
    benchmarked to hang on $n {\geq} 250\mathrm{K}$.)

  \item[3.] \emph{Kasai LCP.} Compute $\mathrm{LCP}[i]$ = length of
    the longest common prefix of $\mathrm{SA}[i{-}1]$ and
    $\mathrm{SA}[i]$ in linear time via the Kasai
    algorithm~\cite{kasai-lcp}.

  \item[4.] \emph{Monotonic-stack enumeration of maximal repeats.}
    Sweep $\mathrm{LCP}$ left-to-right with a stack of
    $(\text{height}, \text{left\_index})$ pairs. For every entry that
    pops because the current LCP dips below it, emit
    $\mathrm{SA}[\text{left\_index}{-}1\,..\,i{-}1]$ as a candidate of
    length \text{height}: those suffixes share at least
    $\text{height}$ characters. This enumerates the nested case
    (a 350-token repeat at freq=3 nested inside a 300-token repeat
    at freq=5 emits both).

  \item[5.] \emph{Floor + sentinel filters.} Discard candidates with
    length $< \ell_{\min}$, with fewer than $f_{\min}$ distinct
    source streams (computed via \texttt{stream\_of\_pos}), or
    containing any sentinel token (negative id). Deduplicate by
    fingerprint hash.

  \item[6.] \emph{Exact-prefix dedup.} Sort surviving candidates by
    length descending. If candidate $A$'s tokens are a prefix of an
    already-kept candidate $B$ with $|A| < |B|$ and they share at
    least one source stream, drop $A$ ($B$ is strictly more
    informative).

  \item[7.] \emph{MinHash near-dup dedup.} Compute a $P$-permutation
    MinHash signature over $k$-shingles for each remaining span. Two
    spans whose signature Jaccard $\geq J$ are merged, keeping the
    longer representative. We use $P{=}64$, $k{=}5$, $J{=}0.8$; the
    blake2b-keyed hash family makes the signatures deterministic
    across runs and tractable on 64-bit machines.

  \item[8.] \emph{Sort + return.} Sort by
    $(-\text{frequency}, -\text{length})$ so callers can take the
    top-$N$ under a budget cap.
\end{enumerate}
\caption{\texttt{mine\_spans}: extract length-$\geq{}\ell_{\min}$,
frequency-$\geq{}f_{\min}$ near-deduped repeated sequences from
a multi-stream token corpus.}
\label{alg:miner}
\end{algorithm}

\paragraph{Empirical scale.} On the union of
\dataset's three subsets (411k tokens, 182 streams) with
$\ell_{\min}=256$, $f_{\min}=3$, $J=0.8$, the miner returns
\textbf{181 spans totalling 278k tokens}, in approximately 18~s on
a single MacBook Pro core. The top hit is a 1{,}691-token span
occurring in 26 of the 182 streams (the CC system prompt with tool
list); the long tail of mid-frequency spans corresponds to repeated
skill preambles and tool-output formatting blocks.

\section{Reproducibility}
\label{app:reproducibility}

\paragraph{Pod configuration.} All \S\ref{sec:evaluation} measurements
use a single pod recipe: \mbox{2$\times$~NVIDIA~H100~80GB} on RunPod's
\texttt{SECURE} cloud, vLLM v0.19.0 with the
\texttt{LMCacheConnectorV1} attached, lmcache v0.4.2 with the
patches in \S\ref{sec:patches}, model
\texttt{meta-llama/Llama-3.3-70B-\allowbreak{}Instruct} loaded at
\texttt{dtype=fp8}, tensor-parallel degree~2, and
\texttt{max\_model\_len=32768}. The cacheblend condition additionally
exports the eight \texttt{LMCACHE\_*} environment variables in
Table~\ref{tab:repro-env}; \texttt{PYTHONHASHSEED=0} is mandatory
for chunk-hash stability across the connector's per-role
\texttt{TokenDatabase} inits.

\begin{table}[t]
\small
\centering
\begin{tabular}{ll}
\toprule
Variable & Value \\
\midrule
\texttt{PYTHONHASHSEED} & \texttt{0} (required) \\
\texttt{VLLM\_KV\_CACHE\_TYPE} & \texttt{lmcache} \\
\texttt{LMCACHE\_USE\_EXPERIMENTAL} & \texttt{True} \\
\texttt{LMCACHE\_LOCAL\_CPU} & \texttt{True} \\
\texttt{LMCACHE\_MAX\_LOCAL\_CPU\_SIZE} & \texttt{40} \\
\texttt{LMCACHE\_CHUNK\_SIZE} & \texttt{256} \\
\texttt{LMCACHE\_ENABLE\_BLENDING} & \texttt{True} \\
\texttt{LMCACHE\_USE\_LAYERWISE} & \texttt{True} \\
\texttt{LMCACHE\_BLEND\_SPECIAL\_STR} & \texttt{" \# \# "} \\
\texttt{LMCACHE\_BLEND\_CHECK\_LAYERS} & \texttt{1} \\
\texttt{LMCACHE\_BLEND\_RECOMPUTE\_RATIOS} & \texttt{0.15} (default) \\
\bottomrule
\end{tabular}
\caption{Environment variables for the \texttt{cacheblend} condition.
Set on the pod before \texttt{vllm serve} is launched. Omitting any
of the eight \texttt{LMCACHE\_*} variables causes the engine either
to crash at startup or to silently fall through to plain prefix
caching (no rescue).}
\label{tab:repro-env}
\end{table}

\paragraph{Provisioning resilience.} RunPod's 2$\times$~H100 pool is
capacity-tight enough that synchronous HTTP 500 ``no instances
available'' rejections at \texttt{POST /pods} are routine. We
combine three mitigations: (i) \texttt{CLOUD\_TYPE=SECURE} with
\texttt{WAIT\_TIMEOUT\_S=3600} and
\texttt{PROVISIONING\_TIMEOUT\_S=1800} to let an accepted pod queue
patiently for hardware; (ii) a bash retry loop around the
\texttt{skillcacher-bench} invocation for synchronous create
rejections (20 attempts, 120~s sleep between); (iii) sequential
execution of multi-condition runs so the harness cannot multiply
its own capacity demand. The
\texttt{scripts/dev/recompute\_probe.sh} driver in our repository is
the canonical example.

\paragraph{Canonical commands.} The complete SWE-V+\texttt{post\_compact} subset replay
(7~captures, 54~pairs across three conditions) is one invocation:
\begin{lstlisting}[basicstyle=\ttfamily\scriptsize, frame=single,
                   framesep=4pt, breaklines=true]
SKILLCACHER_BACKEND_MAX_COMPLETION_TOKENS=512 \
MODEL_NAME=meta-llama/Llama-3.3-70B-Instruct \
DTYPE=fp8 GPU_COUNT=2 TENSOR_PARALLEL_SIZE=2 \
MAX_MODEL_LEN=32768 CLOUD_TYPE=SECURE \
WAIT_TIMEOUT_S=3600 PROVISIONING_TIMEOUT_S=1800 \
skillcacher-bench quality-eval \
  --corpus-dir tests/fixtures/claude_code_real \
  --capture-filter \
"swebench_verified/{astropy-13398,matplotlib-26113,\
pylint-7080,pytest-7490,sphinx-11510},\
post_compact/plan4_compaction_spike_v{6_align,7_norm}" \
  --conditions no_cache,prefix_cache,cacheblend \
  --samples 1 --temperature 0 \
  --out benchmark/results/plan5_quality_section1
\end{lstlisting}
The judge analysis is then a single offline invocation reading
the parquets emitted by the above:
\begin{lstlisting}[basicstyle=\ttfamily\scriptsize, frame=single,
                   framesep=4pt, breaklines=true]
export $(grep -E "^ANTHROPIC_API_KEY=" .env | xargs)
.venv/bin/python scripts/run_judge.py --skip-identical
\end{lstlisting}

\paragraph{Artifact availability.} Three artifacts support
end-to-end reproduction. (i)~\dataset{} on Hugging
Face~\cite{claudecode-trace-hf}, available now under CC-BY 4.0; the
\texttt{scripts/download\_claudecode\_trace.py} wrapper accepts
\texttt{--subset} filters. (ii)~\sys{} source, including the
patched lmcache wheel build script, the proxy, the bench harness,
and the figure-generation scripts; the analysis drivers used for
\S\ref{sec:eval-hitrate} and \S\ref{sec:eval-mode}
(\texttt{scripts/dev/extract\_\allowbreak{}full\_\allowbreak{}corpus\_\allowbreak{}rescue.py},
\texttt{scripts/dev/extract\_\allowbreak{}tool\_\allowbreak{}call\_\allowbreak{}rate.py})
are part of this set. (iii)~The pod-orchestration driver
(\texttt{scripts/dev/oneshot\_pod.\{sh,py\}}) and the recompute-ratio
sweep driver (\texttt{scripts/dev/recompute\_probe.sh}). Items
(ii)~and~(iii) are released under MIT on acceptance, with a frozen
submission-time tag accessible to reviewers via a double-blind-
compatible anonymous Zenodo archive linked from the submission.

\section{Anchor Recognition Walkthrough}
\label{app:anchors}

To make \S\ref{sec:cc-segment}'s parser concrete, this appendix shows
one real CC request body (lightly redacted) with the structural
anchors color-highlighted and the cacheblend separator
\texttt{\textquotedbl{}~\#~\#~\textquotedbl{}} marked at each
injection site.

\definecolor{anchorHeader}{HTML}{FFF3B0}   % header (normalized fields)
\definecolor{anchorGit}{HTML}{B8E0F5}      % gitStatus
\definecolor{anchorCommits}{HTML}{F8C8DC}  % recent commits
\definecolor{anchorReminder}{HTML}{D4EFCE} % system-reminder

\paragraph{Color key.}
\colorbox{anchorHeader}{\strut\,header\,} (per-turn, normalized
per \S\ref{sec:header-norm});
\colorbox{anchorGit}{\strut\,gitStatus\,};
\colorbox{anchorCommits}{\strut\,commits\,};
\colorbox{anchorReminder}{\strut\,\texttt{<system-reminder>}\,}.
The injected separator is shown as
\fbox{\,\texttt{~\#~\#~}\,} between blocks; cacheblend treats
each separated span as one chunk for STORE / LOOKUP / blend.

\begin{figure}[h]
\begin{flushleft}\small\ttfamily
\fbox{\,\texttt{~\#~\#~}\,}\\[1pt]
\colorbox{anchorHeader}{\strut\,x-anthropic-billing-header:
cc\_version=NORM; cch=NORM;\,}\\
\fbox{\,\texttt{~\#~\#~}\,}\\[1pt]
CWD: /Users/$\langle$user$\rangle$/skillcacher\\
Date: 2026-05-08\\[3pt]
\colorbox{anchorGit}{\strut\,gitStatus: This is the git status\dots\,}\\
\colorbox{anchorGit}{\strut\,Current branch: main\,}\\
\colorbox{anchorGit}{\strut\,Status:\,}\\
\colorbox{anchorGit}{\strut\,M~~benchmark/results/audit/$\dots$\,}\\
\colorbox{anchorGit}{\strut\,??~scripts/dev/cacheblend\_quality.sh\,}\\
\fbox{\,\texttt{~\#~\#~}\,}\\[1pt]
\colorbox{anchorCommits}{\strut\,Recent commits:\,}\\
\colorbox{anchorCommits}{\strut\,f8a0cfa fix(capture): plant\dots\,}\\
\colorbox{anchorCommits}{\strut\,12778ec docs: Plan 3 resume\dots\,}\\
\fbox{\,\texttt{~\#~\#~}\,}\\[1pt]
\colorbox{anchorReminder}{\strut\,$\langle$system-reminder$\rangle$\,}\\
\colorbox{anchorReminder}{\strut\,As you answer the user's
questions\dots\,}\\
\colorbox{anchorReminder}{\strut\,$\langle$/system-reminder$\rangle$\,}\\
\fbox{\,\texttt{~\#~\#~}\,}\\[1pt]
\end{flushleft}
\caption{One real CC request, anchors highlighted. The parser
identifies four blocks: the per-turn billing header (yellow), the
gitStatus block (blue), the Recent-commits block (pink), and the
\texttt{<system-reminder>} wrapper (green). Cacheblend separators
are injected on each side of every recognized block. The yellow
block's \texttt{cc\_version} and \texttt{cch} fields have already
been rewritten to \texttt{NORM} by \S\ref{sec:header-norm}'s
header normalizer, which is why the first KV chunk's hash is
stable across multi-turn sessions.}
\label{fig:anchor-walkthrough}
\Description{Color-highlighted representation of a Claude Code
request body. Four anchor classes are marked: yellow for the
per-turn billing header, blue for the gitStatus block, pink for
Recent commits, and green for the system-reminder wrapper.
Cacheblend separator markers shown as boxed double-pound symbols
between each anchor block.}
\end{figure}

\paragraph{What the parser does at each anchor.}
\begin{enumerate}
\item \textbf{Per-turn header (yellow):} detected by the
  \texttt{x-anthropic-\allowbreak{}billing-\allowbreak{}header:}
  prefix. The \texttt{cc\_version} and \texttt{cch} field values
  are rewritten
  in place to \texttt{NORM} placeholders before chunk hashing, so
  the first KV chunk hash is identical across all turns of one CC
  session (\S\ref{sec:header-norm}).
\item \textbf{gitStatus (blue):} detected by the literal
  \texttt{gitStatus:\textbackslash{}n} prefix. Typically 100--400
  tokens long; spans the working-tree summary CC injects on session
  start. Content is stable for the lifetime of a session.
\item \textbf{Recent commits (pink):} detected by
  \texttt{Recent commits:\textbackslash{}n}. Typically 50--200
  tokens; the most recent commit hashes drift across sessions but
  the structural anchor is stable.
\item \textbf{\texttt{<system-reminder>} wrapper (green):} detected
  by the literal tag pair. Variable-length (30--200 tokens
  typically), often falls below
  \texttt{blend\_min\_tokens}=256 in isolation and is rolled into
  an adjacent block rather than getting its own separator pair ---
  see \S\ref{sec:cc-segment}.
\end{enumerate}

\paragraph{Why this matters.} Without the parser, cacheblend's
default \texttt{blend\_special\_str=\textquotedbl{}~\#~\#~\textquotedbl{}}
never appears in CC outbound traffic --- CC's wire format uses
\texttt{\textbackslash{}n\textbackslash{}n} as its block delimiter,
and the resulting 25--180-token spans fall under the 256-token
floor. The parser's job is to inject the configured separator at
exactly the boundaries that produce content-stable, $\geq$256-token
chunks. The four highlighted classes above are what enables
cacheblend to engage at all on natural CC traffic.
